\chapter{Build against Buy}\label{ch:fm:build-against-buy}

In November 2018 Virtu Financial agreed to buy Investment Technology Group, an agency broker, and in March 2019 it completed the purchase for about
\$1.0 billion in cash, \$30.30 a share. Its stated reasons were to offer its clients a complete suite of agency services, ``including transparent
trading and workflow technology, analytics, and liquidity solutions'', running on Virtu's own infrastructure, and about \$123 million a year of
expense savings within two years. A firm can write the software it needs, license it, or buy the company that wrote it together with its clients
and its people. The last is the most expensive way to acquire software, and sometimes the only way to acquire a business built on it.

\section{What firms build and what they buy}

\begin{definition}[Build against buy]\label{def:fm:build-against-buy:bab}
The \emph{build against buy}\index{build against buy} decision is the choice, for a capability the firm needs, between developing it with its own
engineers, licensing or subscribing to a vendor's product, or a mix (a vendor core with the firm's own extensions), over the capability's whole
life.
\end{definition}

The usual rule is that a firm builds what makes it different and buys what makes it the same. A market maker's quoting engine and risk checks are
its edge, and it builds them (Books 11 and 13); its payroll, email and much of its post-trade processing are the same as everyone's, and it buys
them. The hard cases sit in between: order and execution management systems (Book 10), the security master (Book 7), the machine-learning
platform (Book 12). They are needed by everyone, they touch the firm's edge at the edges, and they are expensive both ways.

\begin{definition}[Total cost of ownership]\label{def:fm:build-against-buy:tco}
The \emph{total cost of ownership}\index{total cost of ownership} of a capability is the present value of every cost of having it over its life:
development or licence fees, integration and customisation, infrastructure, support and maintenance, the engineers' time and what they would
otherwise have built, and the cost of leaving it.
\end{definition}

\section{Total cost of ownership}

The chapter's example is an order-management system over seven years (\cref{lst:fm:build-against-buy:costs}). Building it takes seven engineers
for two years, then three a year to maintain it, plus \$0.2 million a year of infrastructure. Buying it costs a \$1.2 million licence rising 5\% a
year, a \$1.0 million integration fee and two engineers in the first year, and one engineer of support every year. An engineer-year costs \$350\,000
loaded, and the firm adds 30\% for the work the engineer would otherwise have done: \$455\,000. The discount rate is 8\%. All figures are
illustrative; Boehm's \emph{Software Engineering Economics} (1981) is the classic reference on estimating the development and maintenance terms.

\omcode{code/firm/buildbuy/firm_buildbuy.py}{41}{60}{The two cost paths and their net present cost: the build's development years and maintenance
tail, and the licence with its escalator, integration and support.}\label{lst:fm:build-against-buy:costs}

\begin{omfigure}
\begin{tikzpicture}
\begin{axis}[width=10.5cm,height=5cm,ybar,bar width=10pt,xlabel={\small year},ylabel={\small cost (\$ million)},ymin=0,ymax=4,
  xtick={1,2,3,4,5,6,7},tick label style={font=\footnotesize},grid=major,grid style={gray!20},table/col sep=comma,
  legend cell align=left,legend style={font=\scriptsize,at={(0.5,-0.3)},anchor=north,/tikz/every even column/.append style={column sep=8pt}},legend columns=2]
\addplot[fill=omDef!55,draw=omDef] table[x=year,y=build]{figdata/desk/20-build-against-buy/paths.csv};
\addplot[fill=omProp!55,draw=omProp] table[x=year,y=buy]{figdata/desk/20-build-against-buy/paths.csv};
\legend{build,buy}
\end{axis}
\end{tikzpicture}

\omcaption{The yearly costs of the two paths: the build's two development years and flat maintenance tail, and the purchase's integration year and
escalating licence. Engineer-years include the 30\% opportunity cost. Data: \texttt{fm\_buildbuy.paths}.}\label{fig:fm:build-against-buy:paths}
\end{omfigure}

Building costs \$3.19 million in each of the first two years and \$1.57 million a year after (\cref{fig:fm:build-against-buy:paths}); buying costs \$3.57 million in the first year and
then \$1.72 million, rising with the licence to \$2.06 million in year seven. Undiscounted, the build costs \$14.2 million over seven years and the
purchase \$14.9 million. Discounting favours buying, since the build's costs come first: the net present costs are \$11.04 million and \$11.30
million (\cref{fig:fm:build-against-buy:horizon}).

\begin{omfigure}
\begin{tikzpicture}
\begin{axis}[width=10.5cm,height=5.2cm,xlabel={\small horizon (years)},ylabel={\small net present cost (\$ million)},xmin=1,xmax=12,ymin=0,ymax=18,
  tick label style={font=\footnotesize},grid=major,grid style={gray!20},table/col sep=comma,
  legend cell align=left,legend style={font=\scriptsize,at={(0.5,-0.3)},anchor=north,/tikz/every even column/.append style={column sep=8pt}},legend columns=2]
\addplot[omDef,thick,mark=*,mark size=1.2pt] table[x=years,y=build]{figdata/desk/20-build-against-buy/horizon.csv};
\addplot[omProp,thick,dashed,mark=square*,mark size=1.2pt] table[x=years,y=buy]{figdata/desk/20-build-against-buy/horizon.csv};
\draw[gray,dotted] (axis cs:7,0) -- (axis cs:7,18);
\legend{build,buy}
\end{axis}
\end{tikzpicture}

\omcaption{Net present cost of building and of buying the order-management system over horizons of one to twelve years. Buying is cheaper from
year 2 to year 6; from year 7 building is cheaper and the gap widens. Data: \texttt{fm\_buildbuy.by\_horizon}.}\label{fig:fm:build-against-buy:horizon}
\end{omfigure}

\begin{proposition}[The engineer cost that decides]\label{prop:fm:build-against-buy:wage}
If every engineer-year on either path costs $w(1+\kappa)$ and all other costs do not depend on $w$, the difference between the net present costs of
building and buying is linear in $w$, $\Delta(w)=a+bw$, with $b$ the present value of the build's engineer-years less the purchase's, times
$1+\kappa$. When $b>0$, building is cheaper exactly when $w<w^*=-a/b$.
\end{proposition}

\begin{proof}
Each year's cost on each path is a constant plus a number of engineer-years times $w(1+\kappa)$; discounting and summing preserve linearity. The
build uses more engineer-years ($b>0$), so $\Delta$ rises with $w$ and changes sign once, at $-a/b$.
\end{proof}

Over seven years the break-even loaded engineer-year is \$363\,000: at the firm's \$350\,000 building wins, by \$0.26 million of present cost,
and a firm paying more for its engineers should buy. From the seventh year on, building stays cheaper at every longer horizon. The margin is thin,
which is the common case: the decision rarely turns on the base case and often on what could go wrong.

\section{Lock-in and the option to switch}

\begin{definition}[Vendor lock-in, switching cost]\label{def:fm:build-against-buy:lockin}
\emph{Vendor lock-in}\index{vendor lock-in} is the dependence on a vendor that makes leaving it costly: data in its formats, workflows built around
its product, integrations, skills and contracts. The \emph{switching cost}\index{switching cost} is what leaving costs: migration, parallel running,
retraining, termination fees and the risk of disruption.
\end{definition}

Lock-in lets a vendor raise its price at renewal by up to the \omterm{def:fm:build-against-buy:lockin}{switching cost} without losing the client. The chapter models the risk on a binomial
tree: at each renewal the licence may rise 30\% with probability 0.3; switching to another vendor costs \$2.0 million and resets the price. With
the price risk, the expected present cost of buying rises from \$11.30 million to \$13.46 million, \$2.42 million more than building; the option to
switch vendors, exercised when it pays, brings it back to \$12.90 million. The option is worth \$0.56 million, and its value is the reason to keep
\omterm{def:fm:build-against-buy:lockin}{switching costs} low (\cref{lst:fm:build-against-buy:switch}). The best policy is simple: switch once the price has risen twice, by 69\%,
except in the last year, when a single year of savings no longer pays for the move (\cref{fig:fm:build-against-buy:tree}).

\omcode{code/firm/buildbuy/firm_buildbuy.py}{79}{106}{The option to switch vendors: backward induction on a tree of renewal prices, with and
without the right to switch.}\label{lst:fm:build-against-buy:switch}

\begin{omfigure}
\begin{tikzpicture}[font=\footnotesize,>=stealth,n/.style={draw,rounded corners,fill=omDef!12,minimum width=1.2cm,align=center},
  s/.style={draw,rounded corners,fill=omProp!30,minimum width=1.2cm,align=center}]
\node[n] (a) at (0,0) {1.00};
\node[n] (b1) at (3,1.1) {1.30}; \node[n] (b0) at (3,-1.1) {1.00};
\node[s] (c2) at (6,2.2) {1.69\\{\scriptsize switch}}; \node[n] (c1) at (6,0) {1.30}; \node[n] (c0) at (6,-2.2) {1.00};
\draw[->] (a) -- node[above,font=\scriptsize] {0.3} (b1); \draw[->] (a) -- node[below,font=\scriptsize] {0.7} (b0);
\draw[->] (b1) -- (c2); \draw[->] (b1) -- (c1); \draw[->] (b0) -- (c1); \draw[->] (b0) -- (c0);
\node[font=\scriptsize] at (0,-3.0) {year 1}; \node[font=\scriptsize] at (3,-3.0) {year 2}; \node[font=\scriptsize] at (6,-3.0) {year 3};
\node[font=\scriptsize,align=left,anchor=west] at (7.0,2.2) {pay \$2.0m, price\\back to 1.00};
\end{tikzpicture}

\omcaption{The first three years of the renewal-price tree: the vendor's price multiplier rises 30\% with probability 0.3 at each renewal. The best
policy switches vendor once the price has risen twice, from year 3 to year 6; in year 7 only after four rises. Data:
\texttt{firm.buildbuy.switch\_policy}.}\label{fig:fm:build-against-buy:tree}
\end{omfigure}

\begin{definition}[Source-code escrow]\label{def:fm:build-against-buy:escrow}
\emph{Source-code escrow}\index{source-code escrow} is an arrangement under which a vendor deposits the source code of its product with an
independent agent, to be released to the client on defined events such as the vendor's insolvency or its ceasing to support the product.
\end{definition}

Escrow does not remove lock-in; it lowers the cost of the worst case, where the vendor disappears, from a rebuild under pressure to a takeover of
code the firm has never run. Its value depends on whether the deposit is kept current and could be built. Real options in general (Dixit and
Pindyck, \emph{Investment under Uncertainty}, 1994) formalise the same idea: flexibility has a value that a static comparison leaves out, and paying
for it (short contracts, standard data formats, escrow) is sometimes cheaper than the lock-in it prevents.

\section{Tutorial: seven years of an order-management system}\label{tut:fm:build-against-buy:tut}

\begin{tutorial}
\textbf{Goal.} Compare building and buying an order-management system over seven years, value the option to switch vendors, and find which inputs
decide the answer. \textbf{End state:} the cost curves by horizon and the tornado chart (\cref{fig:fm:build-against-buy:tornado}).
\begin{tutsteps}
\item \textbf{The paths.} \texttt{fm\_buildbuy.paths()} lists the yearly costs; \texttt{firm.buildbuy.npc} discounts them.
\item \textbf{The break-evens.} \texttt{breakeven\_wage} solves \cref{prop:fm:build-against-buy:wage}; \texttt{breakeven\_horizon} finds the horizon from
which building stays cheaper.
\item \textbf{The option.} \texttt{switch\_value} values the right to switch vendors on the renewal-price tree.
\item \textbf{The tornado.} \texttt{fm\_buildbuy.tornado()} moves each input by 20\% up and down and records the build-minus-buy difference.
\end{tutsteps}
\end{tutorial}

\begin{omfigure}
\begin{tikzpicture}
\begin{axis}[width=10.5cm,height=5.6cm,xmin=-2,xmax=1.5,xlabel={\small build minus buy, 7-year net present cost (\$ million)},
  ytick={0,1,2,3,4,5,6,7},yticklabels={licence,wage,dev engineers,maint engineers,support engineers,kappa,escalator,integration fee},
  y dir=reverse,ymin=-0.6,ymax=7.6,tick label style={font=\footnotesize},grid=major,grid style={gray!20},table/col sep=comma,
  legend cell align=left,legend style={font=\scriptsize,at={(0.5,-0.25)},anchor=north,/tikz/every even column/.append style={column sep=8pt}},legend columns=2]
\addplot[draw=none,mark=none,error bars/.cd,x dir=both,x explicit,error bar style={line width=5pt,omDef!70},error mark=none]
  table[x=lmid,y expr=\thisrow{pos}-0.13,x error=lhalf]{figdata/desk/20-build-against-buy/tornado.csv};
\addplot[draw=none,mark=none,error bars/.cd,x dir=both,x explicit,error bar style={line width=5pt,omProp!70},error mark=none]
  table[x=hmid,y expr=\thisrow{pos}+0.13,x error=hhalf]{figdata/desk/20-build-against-buy/tornado.csv};
\draw[black,thick] (axis cs:-0.2594,-0.6) -- (axis cs:-0.2594,7.6);
\addlegendimage{omDef!70,line width=5pt}\addlegendimage{omProp!70,line width=5pt}
\legend{,,input 20\% lower,input 20\% higher}
\end{axis}
\end{tikzpicture}

\omcaption{Tornado of the seven-year comparison: the build-minus-buy net present cost when each input moves 20\% down or up; the vertical line is
the base case, $-\$0.26$ million (building cheaper). Negative values favour building. The licence and the engineer wage decide the answer. Data:
\texttt{fm\_buildbuy.tornado}.}\label{fig:fm:build-against-buy:tornado}
\end{omfigure}

The licence and the wage dominate: a licence 20\% cheaper makes buying \$1.17 million cheaper than building, one 20\% dearer makes building \$1.69
million cheaper; the wage moves the answer by the same amounts in the other direction. The number of development engineers comes next. The
escalator and the integration fee barely matter over seven years. The decision therefore rests on two numbers that can be negotiated or measured:
the licence the vendor will accept, and what the firm's engineers really cost and could otherwise build.

\section{Choosing a vendor}

\begin{definition}[Request for proposal]\label{def:fm:build-against-buy:rfp}
A \emph{request for proposal}\index{request for proposal} is a structured invitation to vendors to propose a product or service against the firm's
written requirements, in a common format, so that proposals can be scored and compared on functions, performance, cost, terms and risk.
\end{definition}

\begin{method}[Running a build-against-buy decision]\label{met:fm:build-against-buy:run}
\begin{enumerate}
\item Write the requirements, separating what makes the firm different (build or extend) from what it shares with everyone (buy).
\item Issue a \omterm{def:fm:build-against-buy:rfp}{request for proposal} to at least three vendors; score functions, latency and capacity against the firm's service-level agreements
(Book 14), and price the full cost of ownership, not the licence.
\item Estimate the build honestly, with the maintenance tail and the engineers' opportunity cost.
\item Compare over the capability's expected life, with the price risk and the \omterm{def:fm:build-against-buy:lockin}{switching cost}; value the option to switch and what it costs to keep.
\item Negotiate terms that keep the option: data export, standard formats, escrow, price caps at renewal, termination assistance.
\item Revisit at every renewal: lock-in grows with every integration.
\end{enumerate}
\end{method}

A handoff specification (Book 12) between the firm's systems and the vendor's is the cheapest lock-in insurance: if the interface is the firm's, a
new vendor or an in-house build plugs into the same socket. A \omterm{def:fm:technology-strategy:teams}{platform team} (chapter 19) owns it.

\section{Build: the build-against-buy model}\label{bld:fm:build-against-buy:buildbuy}

\begin{build}
\bfield{Purpose} The cost of building and of buying over a horizon, the break-even engineer cost and horizon, the option to switch vendors, and
the sensitivity of the answer.
\bfield{Interface} \texttt{firm.buildbuy}: \texttt{Build}, \texttt{Buy}, \texttt{build\_costs}, \texttt{buy\_costs}, \texttt{npc},
\texttt{breakeven\_wage}, \texttt{breakeven\_horizon}, \texttt{switch\_value}, \texttt{tornado}.
\bfield{Rules} Costs are paid at year ends; engineer-years carry the opportunity cost on both paths; the option's value is never negative.
\bfield{Acceptance tests} \texttt{code/firm/buildbuy/tests/}: cost paths on hand numbers; the break-even wage equalises the costs; the option value
grows with the price risk; the tornado's order.
\bfield{Stretch} A mixed path (vendor core with in-house extensions); insourcing as the switch; the probability that the build overruns.
\end{build}

\begin{omsources}
\item Virtu Financial, Forms 8-K of 7 November 2018 and 1 March 2019 (press releases on the acquisition of Investment Technology Group).
\item A.~K.~Dixit and R.~S.~Pindyck, \emph{Investment under Uncertainty}, Princeton University Press, 1994.
\item B.~W.~Boehm, \emph{Software Engineering Economics}, Prentice-Hall, 1981.
\end{omsources}

\section{Exercises}

\begin{exercise}[$\star$]\label{exo:fm:build-against-buy:1}
What does a loaded engineer-year cost the firm, including the opportunity cost, and what do the build's first two years cost?
\end{exercise}

\begin{exercise}[$\star$]\label{exo:fm:build-against-buy:2}
What does the licence cost in year seven, and what is the purchase's cost that year?
\end{exercise}

\begin{exercise}[$\star$]\label{exo:fm:build-against-buy:3}
Over which horizons is buying cheaper, and from which horizon is building cheaper?
\end{exercise}

\begin{exercise}[$\star\star$]\label{exo:fm:build-against-buy:4}
Why does discounting favour buying in this example?
\end{exercise}

\begin{exercise}[$\star\star$]\label{exo:fm:build-against-buy:5}
What was Virtu's stated rationale for buying ITG, and what did it pay?
\end{exercise}

\begin{exercise}[$\star\star$]\label{exo:fm:build-against-buy:6}
Why is the option to switch worth more when the vendor's price is riskier? What would make it worthless?
\end{exercise}

\begin{exercise}[$\star\star\star$]\label{exo:fm:build-against-buy:7}
\emph{Coding.} Value the option to switch with a price-rise probability of 0.5 instead of 0.3. What does buying cost in expectation, with and
without the option?
\end{exercise}

\begin{exercise}[$\star\star\star$]\label{exo:fm:build-against-buy:8}
\emph{Find the flaw.} ``Building is free: our engineers are already on the payroll.''
\end{exercise}

\section{Problem: Seven Years of an Order-Management System}

\begin{problem}[{Weekend problem --- seven years of an order-management system}]\label{pb:fm:build-against-buy:1}
A chief technology officer must recommend whether to build or buy the firm's next order-management system.

\medskip\noindent\textbf{Part I --- The decision.}
\begin{enumerate}
\item Define \omterm{def:fm:build-against-buy:bab}{build against buy} and the \omterm{def:fm:build-against-buy:tco}{total cost of ownership}.
\item Which capabilities should a trading firm build, and which buy?
\item What did Virtu buy in 2019, for how much, and why?
\item Why is buying the company the most expensive way to acquire software?
\end{enumerate}

\medskip\noindent\textbf{Part II --- The costs.}
\begin{enumerate}[resume]
\item List the yearly costs of each path.
\item Give the net present costs over seven years and the undiscounted totals.
\item State and prove \cref{prop:fm:build-against-buy:wage}.
\item Give the break-even wage and the break-even horizon.
\end{enumerate}

\medskip\noindent\textbf{Part III --- Lock-in.}
\begin{enumerate}[resume]
\item Define \omterm{def:fm:build-against-buy:lockin}{vendor lock-in} and \omterm{def:fm:build-against-buy:lockin}{switching cost}.
\item Describe the renewal-price tree and give the expected cost of buying with and without the option to switch.
\item Define \omterm{def:fm:build-against-buy:escrow}{source-code escrow} and say what it does and does not do.
\item Which contract terms preserve the option to switch?
\end{enumerate}

\medskip\noindent\textbf{Part IV --- The recommendation.}
\begin{enumerate}[resume]
\item Read the tornado: which inputs decide?
\item What would you negotiate with the vendor?
\item What would you measure about your own engineers?
\item Define a \omterm{def:fm:build-against-buy:rfp}{request for proposal} and outline one for the system.
\item How does a handoff specification reduce lock-in?
\item When should the decision be revisited?
\item State the \emph{named result}: the fully loaded engineer-year cost at which building beats buying over seven years, and the value of the
option to switch.
\item In two sentences, write the recommendation.
\end{enumerate}
\end{problem}

\section{Interview questions}

\begin{interviewq}[$\star$ \iqroles{developer}]\label{iq:fm:build-against-buy:1}
Your team wants to build its own market data handler rather than use a vendor's. What would you ask?
\end{interviewq}

\begin{interviewq}[$\star$ \iqroles{trader}]\label{iq:fm:build-against-buy:2}
Why might a trading firm buy an order-management system but build its execution algorithms?
\end{interviewq}

\begin{interviewq}[$\star\star$ \iqroles{researcher}]\label{iq:fm:build-against-buy:3}
How would you value the option to switch vendors?
\end{interviewq}

\begin{interviewq}[$\star\star$ \iqroles{developer}]\label{iq:fm:build-against-buy:4}
What makes a vendor's product hard to leave, and how would you design around it?
\end{interviewq}

\begin{interviewq}[$\star\star$ \iqroles{risk}]\label{iq:fm:build-against-buy:5}
What risks does a critical vendor add, and how would you manage them?
\end{interviewq}

\begin{interviewq}[$\star\star\star$ \iqroles{developer, researcher}]\label{iq:fm:build-against-buy:6}
Estimate the seven-year \omterm{def:fm:build-against-buy:tco}{total cost of ownership} of building a system that takes six engineers a year to build.
\end{interviewq}
